\subsection{The small-index obstruction and the audit of the cardinal Kan--Thurston argument}

\begin{result}{D.K1}{Proved}{Acyclic maps are epimorphisms in the unpointed homotopy category}
Let $t:X\to Y$ be a map of small CW spaces with nonempty, connected,
integrally acyclic homotopy fibres. For every small space $Z$,
$t^*:[Y,Z]\to[X,Z]$ is injective. More precisely the codiagonal
$c:Y\cup_X^hY\to Y$ is a weak equivalence. This verifies GALOIS B.9
in the exact variance needed for CATHOT E.17; the codiagonal
characterization also appears in
\citet[Theorem~2.9]{ABuchholtzDeJongRijke2025}.
\end{result}
\begin{proof}
For connected $Y$, the fibre sequence implies that $X$ is connected
and $\pi_1t$ is onto. Realize the homotopy pushout by the double
mapping cylinder $P$, with its two inclusions $j_0,j_1:Y\to P$.
Van Kampen identifies $\pi_1P$ with
$\pi_1Y*_{\pi_1X}\pi_1Y$, and the codiagonal induces an isomorphism
to $\pi_1Y$ because both maps out of $\pi_1X$ are the same
surjection. An acyclic map is a homology isomorphism with every
local system on its target
\citep[Definition/Proposition~1.2]{AHausmannHusemoller1979}.
For a local system $L$ on $Y$, compare the Mayer--Vietoris sequences
for $Y\leftarrow X\to Y$ and $Y\leftarrow Y\to Y$. The two outer
maps are identities and the middle map is $t$, so
$H_*(P;c^*L)\to H_*(Y;L)$ is an isomorphism. Taking
$L=\mathbb Z[\pi_1Y]$ and using the established fundamental-group
isomorphism identifies this with the homology comparison of the
universal covers. Homological Whitehead makes $c$ a weak
equivalence.

Since $cj_0=cj_1=\mathrm{id}_Y$, it follows that $[j_0]=[j_1]$.
A homotopy $at\simeq bt$, for $a,b:Y\to Z$, gives an actual map
$h:P\to Z$ with $hj_0=a$ and $hj_1=b$. Hence $[a]=[b]$.
For disconnected spaces the construction and argument apply to
each component; acyclicity gives a bijection on components of
$X$ and $Y$. Maps from a disjoint union of spaces are taken
componentwise, with the homotopies assembled on that disjoint
union. This includes the empty case and makes no interchange of
$\pi_0$ with infinite products of ordinary categories.
\end{proof}

\begin{result}{D.K2}{Proved}{The envelope construction used by GALOIS B.10}
For every small group $H$ and every infinite small cardinal
$\lambda\geq|H|$ there is an embedding $H\hookrightarrow A$,
where $A$ is nonabelian, simple, integrally acyclic, and has
cardinality exactly $\lambda$.
\end{result}
\begin{proof}
We check the two constructions in the supplied proof. First enlarge
$H$ to $H*F_\lambda$. For every pair $g\ne1,h$ of elements of a
current group $L$ of cardinality $\lambda$, add a different free
generator $t=t_{g,h}$ and form
\[
 w=gtg^{-1}t^{-1}gt^2g^{-1}t^{-2}.
\]
The words $w$ and $wh$ have infinite order in $L*\langle t\rangle$.
The first is cyclically reduced with eight alternating syllables;
cyclic reduction of the second has eight syllables unless $hg=1$,
in which case it has six. In all cases its length is at least two
and every displayed free-product syllable is nonidentity. Therefore
the two infinite cyclic subgroups can be identified by an HNN
extension with relation $s^{-1}ws=wh$. The simultaneous HNN
embedding theorem preserves the old group
\citep[Theorems~I--II]{AHigmanNeumannNeumann1949}. The relation puts
$h$ in the normal closure of $g$. Countably repeating this step
gives a simple overgroup of cardinality $\lambda$: any nonidentity
$g$ and any $h$ eventually occur together and satisfy that normal
generation condition. Noncommutativity follows from the embedded
free group.

For the acyclic envelope, start with $L\times L$ and its two
inclusions $j(h)=(h,1)$, $k(h)=(1,h)$. Add stable letters $d,s$
with
\[
 d^{-1}j(h)d=j(h)k(h),\qquad s^{-1}j(h)s=k(h).
\]
The maps identify embedded copies of $L$, so HNN again gives an
embedding of $L\times L$. Call the resulting group $m(L)$.
The homology calculation is the mitosis calculation of
\citet[Proposition~4.1 and Theorem~4.2]{ABaumslagDyerHeller1980}.
Explicitly, let $f:G\to L$ induce zero in positive degrees below
$n$ over a field. In degree $n$ the K\"unneth decomposition of the
diagonal following $f$ has only its two axis terms, since all
mixed positive-degree terms vanish. In $m(L)$ each axis map and
the diagonal map are conjugate to $j$. Hence on $H_n(G)$,
$(jf)_*=(jf)_*+(jf)_*$, so $(jf)_*=0$. Induction shows that a
map from one stage to its $n$-th successor in the iteration of
$m$ kills positive homology through degree $n$. Each bar cycle
uses finitely many group elements. The countable union is
therefore acyclic over every field. Rational coefficients make
its positive integral homology torsion; the universal coefficient
sequence in the next degree with $\mathbb F_p$ coefficients
excludes $p$-torsion for every prime. Its positive integral
homology is zero. The construction does not increase the
infinite cardinal $\lambda$.

Alternate simple and acyclic envelopes, always by embeddings,
and take their countable union $A$. Given $g\ne1$ and $h\in A$,
both lie in a later simple stage, so the normal closure of $g$
in $A$ contains $h$. Each finite integral bar cycle lies in a
later acyclic stage and bounds there. Thus $A$ has all the stated
properties. This verifies the cardinal assertion as well as the
homological assertion, rather than obtaining a merely countable
special case from a citation.
\end{proof}

\begin{result}{D.K3}{Proved}{Verification of the cardinal-controlled Kan--Thurston construction}
For every connected $\mathcal U$-small CW complex $Y$ and every
$\mu\in\mathcal U$ there are a $\mathcal U$-small group $G$ and
an acyclic map $t:BG\to Y$ such that $\pi_1t$ is onto and every
homomorphism $G\to L$ with $|L|\leq\mu$ kills
$P=\ker(\pi_1t)$. In particular the proposed $(\mathrm{KT}_\mu)$
is true with the quantifiers in the prompt. The induction in
GALOIS B.11, including its subcomplex-injectivity step, proves it.
\end{result}
\begin{proof}
We give the checked induction in full, since injectivity is the
material issue in the inherited proof. Choose a small simplicial
complex $K$ and a homotopy equivalence $h:|K|\to Y$
\citep[Theorem~2C.5]{AHatcher2002}. Choose an infinite
$\lambda\in\mathcal U$ exceeding $\mu$ and at least the number
of simplices of $K$. The construction makes, simultaneously for
every subcomplex $L\subseteq K$, a space $T(L)$ and an actual
map $t_L:T(L)\to|L|$. The choices attached to each simplex are
made once and reused in every subcomplex containing it. The
induction assertions are that $T(L)$ is aspherical on each
component and that every inclusion $T(L')\hookrightarrow T(L)$
is injective on fundamental groups, component by component.

In dimensions zero and one put $T(L)=|L|$ and $t_L=\mathrm{id}$.
Reduced edge paths prove both assertions for graphs. Suppose
the construction is made through dimension $n-1$, where
$n\geq2$. For every $n$-simplex $\sigma$ choose by D.K2 an
embedding
\[
 \pi_1T(\partial\sigma)\hookrightarrow C_\sigma
\]
into a simple acyclic group of cardinality $\lambda$.
The connected aspherical source admits a cellular map
$a_\sigma:T(\partial\sigma)\to BC_\sigma$ realizing this
embedding. Let $Q_\sigma$ be its mapping cylinder, with
$T(\partial\sigma)$ at the free end. Thus this end is a CW
subcomplex, its map to $Q_\sigma$ is injective on $\pi_1$,
and $Q_\sigma\simeq BC_\sigma$ is aspherical and acyclic.
Define
\[
 T(L)=T(L^{(n-1)})\cup_{\coprod_{\sigma\in L_n}
                   T(\partial\sigma)}
                   \coprod_{\sigma\in L_n}Q_\sigma.
 \tag{D.K3.1}
\]
On a cylinder, $t_L$ moves $t_{\partial\sigma}(x)$ linearly
to the barycentre of $\sigma$, and it sends $BC_\sigma$ to
that barycentre. This defines the required continuous map,
compatible with all subcomplex inclusions.

We use the following graph-of-spaces fact with its precise
injectivity hypothesis. A graph of connected aspherical CW
vertex spaces, with aspherical edge spaces attached by mapping
cylinders and injective edge homomorphisms, is aspherical and
embeds its vertex groups. The amalgamated-product and HNN
normal forms give the embeddings. In its universal cover,
the incidence graph is the tree of cosets. Every vertex and
edge inverse image is contractible. The space over any finite
subtree is contractible by successively attaching a
contractible vertex along a contractible edge cylinder. A
sphere map and any finite cellular data lie over a finite
subtree, so the universal cover is weakly contractible and
hence contractible. This is also the amalgamation argument
recalled in \citet[Section~1, before Proposition~1.1]
{ABaumslagDyerHeller1980}. Components of a disconnected edge
space are treated as separate edges.

In (D.K3.1) each boundary group injects into the lower-skeleton
group by induction and into $C_\sigma$ by construction. Thus
$T(L)$ is aspherical, and the lower skeleton and each
$BC_\sigma$ embed on $\pi_1$. For the stronger subcomplex
assertion, fix $L\subseteq K'$ of dimension at most $n$ and
form
\[
 V=T((K')^{(n-1)})
       \cup_{T(L^{(n-1)})}T(L).
 \tag{D.K3.2}
\]
This is a graph of spaces with one edge for each component
of $T(L^{(n-1)})$. Its two edge inclusions are injective:
the first by the dimension-$n-1$ induction hypothesis, the
second by the preceding application to $T(L)$. Therefore
$T(L)\to V$ and $T((K')^{(n-1)})\to V$ are injective on
fundamental groups. Attach the remaining $Q_\sigma$ for
$\sigma\in K'_n\setminus L_n$. Their boundary groups inject
through $T((K')^{(n-1)})$ and keep doing so after each
attachment. The result is exactly $T(K')$, by the same
pushout construction and the choices fixed simplexwise.
Hence $\pi_1T(L)\to\pi_1T(K')$ is injective. This proves
the required step without treating an arbitrary map of
aspherical spaces as injective. Filtered unions over
dimensions preserve these assertions because spheres,
nullhomotopies, and group relations use finitely many cells.

For every local system $M$ of abelian groups on $|L|$ we
claim
\[
 H_*(T(L);t_L^*M)\xrightarrow{\cong}H_*(|L|;M).
 \tag{D.K3.3}
\]
It holds for graphs. In an $n$-simplex the coefficient
system is constant. The map of pairs
$(Q_\sigma,T(\partial\sigma))\to(\sigma,\partial\sigma)$
is an isomorphism on relative homology by the long exact
sequences, acyclicity of $Q_\sigma$, and induction for the
boundary. Excision gives the direct sum of these relative
comparisons over all $n$-simplices of $L$. The skeletal
long exact sequence proves (D.K3.3) at that dimension;
cellular chains with local coefficients commute with the
union over skeleta. Thus it holds without a dimension
bound. All spaces constructed have at most $\lambda$
cells: bar models of the groups do, and the finite products
and countable unions involved preserve that bound.

Put $G=\pi_1T(K)$ and let $N$ be the normal closure in $G$
of all the embedded $C_\sigma$, $\dim\sigma\geq2$.
The map $t_K$ kills these groups. Its induced quotient map
is the isomorphism
\[
 G/N\xrightarrow{\cong}\pi_1|K|.
 \tag{D.K3.4}
\]
Indeed, at dimension two, killing $C_\sigma$ kills exactly
the boundary relator of that simplex. For $n\geq3$,
the inductive quotient assertion for $\partial\sigma$
and $\pi_1\partial\sigma=0$ say that the boundary group
is normally generated by the already attached lower
cone groups. It introduces no new relation after those
groups have been killed. Van Kampen proves (D.K3.4) by
induction and by the union over dimensions. Each
$C_\sigma$ is perfect, so their normal closure $N$ is
perfect: its generating conjugate subgroups lie in
$[N,N]$.

Choose a based homotopy equivalence $b:BG\to T(K)$ and
set $t=h t_K b$. Equation (D.K3.3), with
\citet[Definition/Proposition~1.2]{AHausmannHusemoller1979},
makes this specified map acyclic. Equation (D.K3.4)
identifies its perfect kernel with $N$. A homomorphism
$G\to L$, $|L|\leq\mu$, restricts trivially to each
$C_\sigma$: simplicity makes a nontrivial restriction
injective, which its cardinality $\lambda>\mu$ forbids.
It therefore kills their normal closure. When $K$ has
dimension at most one there are no cone groups and the
kernel is trivial, so that case is included.
\end{proof}

\begin{result}{D.K4}{Proved}{The small-index version of CATHOT E.17}
Let $\mathcal H=\mathrm{Cat\,hot}_{\mathcal U}$, whose morphisms
are functors modulo finite zigzags of natural transformations.
Let $I$ be a nonempty $\mathcal U$-small filtered category and
$K:I\to\mathcal H$. If
\[
 H_K(X)=\operatorname{colim}_{i\in I}
                 \operatorname{Hom}_{\mathcal H}(X,K_i)
\]
inverts $W_\infty$ on all $\mathcal U$-small categories $X$,
then it inverts $W_1$. Moreover the specified unit
\[
 (K_i)_{i\in I}\longrightarrow(\Pi_1K_i)_{i\in I}
\]
is an isomorphism in $\operatorname{Ind}_{\mathcal U}(\mathcal H)$,
without changing the indexing category. This verifies GALOIS B.12
and settles the numbered target CATHOT E.17 in that small-index scope.
\end{result}
\begin{proof}
The functor $q:\mathcal H\to\mathrm{Hot}$ is localization at
nerve weak equivalences. A natural transformation supplies a
functor on $X\times[1]$; the two endpoint inclusions become equal
after its projection to $X$ is inverted. Thus passing from
$\mathrm{Cat}$ to $\mathcal H$ before localization has no effect.
The assumed $H_K$ is therefore the pullback of a presheaf $H'$
on $\mathrm{Hot}$, and its natural comparison is
\[
 H'(Y)\cong\operatorname{colim}_{i\in I}[Y,BK_i].
 \tag{D.K4.1}
\]
To justify this crucial step without assuming each $K_i$ local,
take presheaves in a larger universe. Restriction $q^*$ is fully
faithful. Its left adjoint, left Kan extension along
$q^{\mathrm{op}}$, sends a representable to the corresponding
representable on $\mathrm{Hot}$ by Yoneda, and preserves colimits.
Its unit is invertible on the presheaf already in the image of
$q^*$. Applying it to the displayed colimit proves (D.K4.1),
including its naturality and its specified comparison map.

Choose a small infinite cardinal $\mu$ bounding the numbers of
arrows of every $K_i$. Such a uniform bound exists because $I$
and the family are small in one universe. For every component
$Y_c$ of a small CW space $Y$ apply D.K3 to obtain
$t_c:BG_c\to Y_c$. Put $T=\coprod_c BG_c$, $t=\coprod_c t_c$,
and let $v:Y\to\tau_{\leq1}Y=\coprod_c B\pi_1Y_c$ be first
truncation in the homotopy category. The quotient homomorphisms
$G_c\twoheadrightarrow\pi_1Y_c$ represent $vt$. Write $\mathbf B G$
for the one-object groupoid of a group $G$, so that $BG$ denotes
its classifying space. For every $i$
restriction induces an actual isomorphism of ordinary categories
\[
 \operatorname{Fun}\left(\coprod_c \mathbf B\pi_1Y_c,K_i\right)
     \xrightarrow{\cong}
 \operatorname{Fun}\left(\coprod_c \mathbf B G_c,K_i\right).
 \tag{D.K4.2}
\]
A functor on $\mathbf B G_c$ is an object $k$ together with a homomorphism
$G_c\to\operatorname{Aut}_{K_i}(k)$. Its target group has size at
most $\mu$, so it factors uniquely through $\pi_1Y_c$. A natural
transformation is an arrow commuting with these actions; because
the quotient homomorphism is surjective, this is exactly the same
condition after factorization. Thus (D.K4.2) is an isomorphism on
objects and arrows, even when $K_i$ has noninvertible
endomorphisms and when there are infinitely many components.

By (D.K4.1), the resulting bijection is the composite
\[
 H'(\tau_{\leq1}Y)\xrightarrow{H'(v)}H'(Y)
                      \xrightarrow{H'(t)}H'(T).
\]
The second arrow is injective by D.K1: at each stage
$[Y,BK_i]\to[T,BK_i]$ is injective, and a filtered colimit of
these compatible injections remains injective. The composite
is bijective by (D.K4.2), so the first arrow is bijective.
Here taking components was done after an isomorphism of functor
categories; no invalid commutation with an infinite product is
involved. For $Y=BX$ this says exactly that
$H_K(\Pi_1X)\to H_K(X)$, induced by the category groupoid-completion
unit, is a bijection. Since $\Pi_1$ takes $W_1$ to equivalences
of groupoids, $H_K$ inverts $W_1$.

Finally, groupoid completion reflects $\mathcal H$ into the
full subcategory of groupoids modulo natural isomorphisms.
The ind-Hom formula
\[
 \operatorname{Hom}_{\operatorname{Ind}(\mathcal H)}
       ((A_i),(B_j))
   =\lim_i\operatorname{colim}_j
            \operatorname{Hom}_{\mathcal H}(A_i,B_j)
\]
extends that reflection to ind-objects. Write $L K=(\Pi_1K_i)$
and $\eta:K\to LK$. Taking the limit over $i$ of the
bijections $H_K(\Pi_1K_i)\to H_K(K_i)$ supplies
$r:LK\to K$ with $r\eta=\mathrm{id}_K$. The object $LK$ has
the same orthogonality to every $X\to\Pi_1X$, since each
groupoid target does. Precomposition with $\eta$ is therefore
injective on endomorphisms of $LK$. The identity
$(\eta r)\eta=\eta=\mathrm{id}_{LK}\eta$ gives
$\eta r=\mathrm{id}_{LK}$, proving that the specified unit
is an ind-isomorphism. The empty test category is included:
both sides of (D.K4.2) are then terminal.
\end{proof}

\begin{result}{D.K5}{Proved}{Application to V ms.44--46 and the universe restriction}
For a small relative category $(M,\Sigma)$ and fixed objects $a,b$,
the countable path system of V ms.~44--45 is in the scope of D.K4.
Consequently, if
\[
 X\longmapsto
 \operatorname{Hom}_{M^X[\Sigma_X^{-1}]}(a_X,b_X)
\]
inverts weak equivalences of all small $X$ in that same universe,
it factors through $\Pi_1$. This conclusion does not require that
the functor first be represented by a space. If it is represented
by a space $Z$, then $Z$ is a $1$-type. No such conclusion follows
from invariance on a class of shapes bounded relative to a larger
$M$.
\end{result}
\begin{proof}
Every localization morphism in $M^X$ is represented by a finite
word of arrows and formal inverses of pointwise $\Sigma$-arrows.
A single word has one finite length and type for the whole
diagram $X$, however large $X$ is. Padding by identities puts
it at a stage $K_n$ of the countable alternating path system.
Commutative ladders with identity endpoint components give
the same localization morphism: after inverting the backward
arrows, the commutative-square equations telescope from the
first endpoint to the second. No vertical arrow is inverted in
this calculation, so arbitrary vertical components are allowed.
Conversely, the relations
defining localization, namely composition, identities, and
inverting a $\Sigma$-arrow, are represented by finite
reparametrizations and such ladders. A finite proof of equality
uses a common later stage. Applying this description in the
functor category gives the natural bijection of sets
\[
 \operatorname{Hom}_{M^X[\Sigma_X^{-1}]}(a_X,b_X)
       \cong\operatorname{colim}_{n\in\mathbb N}
                      \pi_0\operatorname{Fun}(X,K_n).
\]
This is exactly the variance and countable indexing printed
at V ms.~45. It can also be obtained by applying C.P3 to the
relative category $(M^X,\Sigma_X)$, taking connected components,
and using the literal exponential identification of its
fixed-type path category with $\operatorname{Fun}(X,K_n)$.
Each $K_n$ is small because $M$ is small, so
D.K4 applies. Neither two-out-of-three nor saturation is
needed for this formal localization description.

If $H(X)=[BX,Z]$, its first-truncation invariance implies that
$[\tau_{\leq1}Y,Z]\to[Y,Z]$ is bijective for every small
space $Y$, by categorical models of spaces. At $Y=Z$ the
identity factors through $Z\to\tau_{\leq1}Z$. Thus $Z$ is a
retract in the homotopy category of a $1$-type, and its
homotopy groups above degree one vanish. In a rectified
diagram theory the constant-to-constant formula identifies
$Z$ with the relevant mapping space, giving the stated
restriction on the unbounded small presentation.

The proof chooses groups of size strictly larger than a
uniform bound on the stages. Those classifying categories
are still small in the original universe, but they need
not belong to a prescribed bounded class of admissible
shapes. A model category large relative to those shapes
does not supply the uniform bound used in D.K4. This is
the precise point at which the bounded formulation avoids
the obstruction. The expression ``Thomason--Kan'' at
V ms.~46 can accordingly be read as localness in
$\mathrm{Cat\,hot}$, realized here by groupoid stages;
Thomason fibrancy alone does not justify that reading.
\end{proof}
